Sources : études citées, présentées dans le complément « Les élasticités estimées sur données françaises » ; chaque chiffre, avec sa page et son tableau, figure dans donnees/elasticites\_france.csv.\par Lecture : chaque panneau compare deux estimations ou plus d'une même étude, obtenues sur les mêmes données avec la même méthode ; seule change la dimension que nomme son titre. Les panneaux ne se comparent pas entre eux, et leurs échelles diffèrent. Une élasticité de 0,2 signifie +2 \% de revenu déclaré pour +10 \% de part gardée sur l'euro marginal. Cabannes et al. et Bach et al. ne publient pas d'écart-type ; chez les premiers, seul le décile supérieur est significatif. Chez Lefebvre et al., les pensions, qui ne réagissent pas, sont comptées avec les revenus d'activité.
